\documentclass{standalone}
\usepackage{ctex}
\usepackage{tikzplot}

% Styles
\tikzstyle{axes}=[]
\tikzstyle{important line}=[very thick]
\tikzstyle{information text}=[rounded corners,fill=red!10,inner sep=1ex]

\begin{document}

\begin{tikzpicture}
  \tikzmath{
    % 5 点的齐次坐标
    \Pa{1,1} = 0; \Pa{2,1} = 5; \Pa{3,1} = 1; 
    \Pb{1,1} = -2; \Pb{2,1} = -2; \Pb{3,1} = 1; 
    \Pc{1,1} = 3; \Pc{2,1} = 3; \Pc{3,1} = 1; 
    \Pd{1,1} = -3; \Pd{2,1} = 3; \Pd{3,1} = 1; 
    \Pe{1,1} = 4; \Pe{2,1} = -2; \Pe{3,1} = 1; 
  }
  \tikzset{
    conics/define/five points={\Pa,\Pb,\Pc,\Pd,\Pe,\CCa},
    conics/cross={\Pa,\Pb,\La},
    conics/cross={\Pb,\Pc,\Lb},
    conics/cross={\Pc,\Pd,\Lc},
    conics/cross={\Pd,\Pe,\Ld},
    conics/cross={\Pe,\Pa,\Le},
    conics/define/five tangents={\La,\Lb,\Lc,\Ld,\Le,\CCb},
    mat/normalize={\CCa,3,3,\Ca},
    mat/normalize={\CCb,3,3,\Cb},
    mat/stdout={\Ca,3,3},
    mat/stdout={\Cb,3,3},
    % 构造起始点
    conics/midpoint={\Pb,\Pc,\Pm},
    conics/intersect cl={\Ca,\Pa,\Pm,\Qa,\Qb},
    conics/exclude={\Qa,\Qb,\Pa,\Pstart},
  }

  \tikzmath{
    % 起始点, 每次都修改, 用于排除已找到顶点
    % \PUprev = \Pe{1,1};
    % \PVprev = \Pe{2,1};
    % \PWprev = \Pe{3,1};
    \PUprev = \Pstart{1,1};
    \PVprev = \Pstart{2,1};
    \PWprev = \Pstart{3,1};
    % 多边形与内切圆的切点, 每次都修改, 用于排除已找到的边
    \TUprev = 0;
    \TVprev = 0;
    \TWprev = 1;
  }

  % \clip (-6,-6) rectangle (6,6);
  \axes[grid] (-5:5,-5:5);
  \conic[thick,draw=teal] (\Ca);
  \conic[thick,draw=navy] (\Cb);

  \foreach \i in {1,...,5}{%sandbox
    \tikzmath{
      % 当前的顶点和切点
      \Pcurr{1,1}=\PUprev; \Pcurr{2,1}=\PVprev; \Pcurr{3,1}=\PWprev;
      \Tcurr{1,1}=\TUprev; \Tcurr{2,1}=\TVprev; \Tcurr{3,1}=\TWprev;
    }
    
    \tikzset{
      mat/stdout={\Pcurr,3,1},
      mat/stdout={\Tcurr,3,1},
      conics/tangents={\Cb,\Pcurr,\Ta,\Tb},
      conics/dehomogenize={\Pcurr,\PPa},
      conics/exclude={\Ta,\Tb,\Tcurr,\Tnext},
      conics/intersect cl={\Ca,\Pcurr,\Tnext,\P,\Q},
      conics/exclude={\P,\Q,\Pcurr,\Pnext},
      conics/dehomogenize={\Pnext,\PPn},
      conics/dehomogenize={\Tnext,\TTn},
    }
    % % 导出下一个顶点和切点
    \xdef\PUprev{\Pnext{1,1}}
    \xdef\PVprev{\Pnext{2,1}}
    \xdef\PWprev{\Pnext{3,1}}
    \xdef\TUprev{\Tnext{1,1}}
    \xdef\TVprev{\Tnext{2,1}}
    \xdef\TWprev{\Tnext{3,1}}
 
    \coordinate (P) at (\PPa{1},\PPa{2});
    \coordinate (Q) at (\PPn{1},\PPn{2});
    \coordinate (T) at (\TTn{1},\TTn{2});
    \draw[red] (P) -- (Q);
    \foreach \p/\placement in {P/above,T/above}{
      \fill[red] (\p) circle (1.5pt);
     \draw (\p) node[\placement] {$\p\i$};
    }
  }

  \draw[xshift=-6.15cm,yshift=-6.15cm] node [right,text width=12cm,style=information text]
    {
      \textsf{{\color{red}{彭赛列闭合定理 Poncelet's Porism}} 若存在一个封闭多边形外切其中一条圆锥曲线且内接另一条圆锥曲线，则此封闭多边形内接的圆锥曲线上每一个点都是满足这样（切、接）性质的封闭多边形的顶点，并且所有满足该性质的封闭多边形的边数相同。}
    };
\end{tikzpicture}

\end{document}